\newpage
% *************************
% **   FACTIBILIDAD      **
% *************************
\section{Estudio de Factibilidad}

El estudio de factibilidad determina si el proyecto puede ejecutarse con éxito dentro de las restricciones reales de tiempo, hardware, presupuesto, datos y acceso institucional. Se estructura en tres dimensiones: viabilidad técnica, viabilidad económica y viabilidad operativa.

La conclusión anticipada es que el proyecto \textbf{sí es factible como prototipo académico} (\textit{Minimum Viable Product}), pero \textbf{no} como sistema de producción ni como modelo entrenado desde cero. La razón es que el plan reduce correctamente el riesgo al proponer \textbf{seleccionar, cuantizar y desplegar} un SLM preentrenado con RAG local y detección de idioma, en lugar de entrenar un LLM desde cero. Esta estrategia es coherente con la literatura reciente sobre despliegue \textit{on-device} \citep{Xu2024} y educación \textit{offline} en entornos de baja conectividad \citep{Walusimbi2026}. La calificación global del proyecto, con la información disponible, es \textbf{AMARILLO favorable}: el mayor riesgo no es el modelo base, sino la curación de datos, la pertinencia lingüística y la validación pedagógica intercultural.

\subsection{Factibilidad técnica y de datos}

La viabilidad técnica es \textbf{realista dentro de un perímetro bien acotado}. El plan propone un asistente textual \textit{offline} con detección de idioma, enrutamiento lingüístico, RAG local y un SLM cuantizado; excluye voz y sitúa la contribución en la integración y evaluación del \textit{pipeline} completo, no en el entrenamiento de un modelo nuevo. Esa acotación reduce el riesgo de sobreingeniería y hace el proyecto ejecutable dentro del plazo de una tesis.

\subsubsection{Hardware e inferencia local}

El archivo GGUF de \textbf{Qwen2.5-7B-Instruct en cuantización Q4\_K\_M pesa 4,68 GB}; el de \textbf{Qwen2.5-3B-Instruct pesa 2,1 GB}. El consumo real de RAM en ejecución es mayor que el tamaño del archivo, porque incluye los pesos, la caché KV y el \textit{overhead} del \textit{runtime}: aproximadamente \textbf{6–7 GB} para el modelo de 7B y \textbf{3 GB} para el de 3B. Esto implica que en una laptop con \textbf{16 GB de RAM}, el 7B corre con holgura; en una con \textbf{8 GB}, el 7B es marginal y el \textbf{3B es la opción segura}. A efectos de este proyecto, se adopta el \textbf{Qwen2.5-3B como configuración base} (equipos de 8 GB) y el \textbf{Qwen2.5-7B como variante de calidad} (equipos de 16 GB).

La inferencia en CPU está limitada por el ancho de banda de memoria, no por la frecuencia del procesador. En una laptop de gama media (Intel Core i5 de 10.ª a 13.ª generación o AMD Ryzen 5, con memoria DDR4/DDR5), un 7B Q4\_K\_M genera aproximadamente \textbf{8–12 tokens por segundo} y un 3B genera \textbf{10–25 tokens por segundo}. Estas velocidades son aceptables para un asistente educativo de turnos cortos. La cuantización no exige GPU; se realiza con utilidades de \texttt{llama.cpp} en CPU o, para acelerar la conversión, en el nivel gratuito de Google Colab (GPU T4) o Kaggle, sin costo adicional.

\subsubsection{Datos y corpus disponible}

El corpus paralelo Quechua Collao–español de \citet{Orcotoma2025} está confirmado: \textbf{44 263 palabras} digitalizadas por OCR a partir de cinco libros de la Academia Mayor de la Lengua Quechua (AMLQ). Es un conteo de palabras, no de pares de oraciones; el propio artículo advierte que no se aplicó limpieza \textit{post-OCR}. Como recursos complementarios de dominio público se identifican: el corpus monolingüe \texttt{Llamacha/monolingual-quechua-iic} (base del modelo QuBERT para quechua sureño), datos de AmericasNLP 2023 —que comprenden aproximadamente 154 000 pares quechua–español de entrenamiento— y materiales curriculares EIB abiertos del MINEDU.

\textbf{Advertencia dialectal crítica:} la variante \textit{quy} utilizada en AmericasNLP corresponde al Quechua Chanka (Ayacucho), no al Quechua Collao (Cusco–Puno). Su uso en el RAG exige validación y normalización ortográfica previa por un especialista. Construir la base RAG en doce meses es \textbf{realista} porque el RAG no entrena el modelo; solo requiere curar, segmentar (\textit{chunking}) y vectorizar documentos. El \textit{tokenization premium} del quechua —lengua aglutinante con alta fertilidad de tokens— debe monitorearse como métrica de control del sistema.

El carácter multilingüe de Qwen2.5 no debe confundirse con competencia real en Quechua Collao: el modelo card declara soporte para español y más de 29 idiomas, pero \textbf{no garantiza soporte específico de Quechua Collao} \citep{Court2024}. La pertinencia en esta variedad dependerá del corpus recuperado, los glosarios curados y los \textit{prompts} diseñados con revisión de hablantes nativos, no del conocimiento paramétrico del modelo.

\subsubsection{Stack tecnológico}

Todo el \textit{stack} propuesto es abierto, local y gratuito: \textbf{\texttt{llama.cpp} / Ollama} (inferencia \textit{offline}), \textbf{LangChain o LlamaIndex} (orquestación RAG), \textbf{sentence-transformers} (\textit{embeddings} locales) y \textbf{FAISS o ChromaDB} (base vectorial local). La curva de aprendizaje es moderada y los ecosistemas están ampliamente documentados. El costo recurrente de inferencia es \textbf{S/ 0}.

\subsubsection{Matriz técnica de factibilidad}

\begin{table}[H]
\centering
\caption{Matriz de viabilidad técnica aplicada al proyecto}
\label{tab:factibilidad_tecnica}
\small
\begin{tabular}{@{}p{4.2cm}p{2.5cm}p{4.5cm}p{3.5cm}@{}}
\toprule
\textbf{Componente crítico} & \textbf{Estado actual} & \textbf{Justificación} & \textbf{Contingencia} \\
\midrule
Cómputo para inferencia (CPU-only) & Por gestionar & Inferencia local validada en literatura \citep{Walusimbi2026}; ~8–25 tok/s en laptop gama media según tamaño de modelo & Iniciar con 3B; escalar a 7B solo si el equipo lo soporta \\
\midrule
Dataset y corpus para RAG & Por gestionar & Corpus base disponible \citep{Orcotoma2025}; pequeño, requiere limpieza y complemento curricular & Acotar dominio a un área y grado; curación manual validada por especialista AMLQ \\
\midrule
Stack tecnológico & Parcialmente cubierto & \textit{Stack} abierto y documentado; soporte Qwen en ecosistemas estándar \citep{Xu2024} & Congelar versiones con \texttt{requirements.txt}; evitar dependencias de pago \\
\midrule
Riesgo de sobreingeniería & Alto si no se contiene & Viable como prototipo textual; riesgo si se agrega voz, app móvil o despliegue institucional masivo \citep{Hevia2025} & Interfaz mínima (CLI/web local); priorizar evaluación del \textit{pipeline} \\
\bottomrule
\end{tabular}
\end{table}

\subsection{Factibilidad económica e infraestructura}

Económicamente, el proyecto es \textbf{más viable de lo que podría estimarse inicialmente}, siempre que se mantenga la estrategia \textit{offline}. El presupuesto original del plan asciende a S/ 7 205,00. Sin embargo, el análisis de mercado identifica \textbf{dos partidas mal calibradas}: el cómputo en la nube está sobreestimado (puede reducirse a casi cero con niveles gratuitos) y el almacenamiento externo está subestimado frente al precio real de un SSD en el mercado peruano.

\subsubsection{Fuentes de financiamiento y niveles gratuitos}

Los pasos de cuantización y vectorización son tareas \textbf{puntuales} (\textit{one-off}), no servicios persistentes. Se cubren íntegramente con niveles gratuitos:

\begin{itemize}
    \item \textbf{Kaggle:} hasta 30 horas semanales de GPU con sesiones de hasta 9 horas; suficiente para cuantización y vectorización.
    \item \textbf{Google Colab (nivel gratuito, GPU T4):} adecuado para las mismas tareas; no imprescindible contratar Colab Pro.
    \item \textbf{GitHub Student Developer Pack:} incluye créditos de DigitalOcean, Azure y acceso a Copilot y Codespaces. \textbf{Advertencia:} desde junio de 2026, los créditos de DigitalOcean \textit{no pueden usarse en GPU Droplets}; son útiles para hosting en CPU, no para cómputo GPU. Kaggle/Colab son la vía recomendada.
    \item \textbf{AWS Educate:} formación autoguiada y laboratorios sin costo; no garantiza GPU sostenida para el proyecto.
\end{itemize}

Las APIs de terceros (OpenAI, Google Cloud) deben evitarse por dos razones: aumentan el costo y, al requerir internet, \textbf{contradicen el propio supuesto de la investigación} (asistente \textit{offline}).

\subsubsection{Presupuesto corregido}

\begin{table}[H]
\centering
\caption{Comparación del presupuesto original vs.\ valores de mercado corregidos (junio 2026)}
\label{tab:presupuesto_corregido}
\small
\begin{tabular}{@{}p{6cm}rrp{3.5cm}@{}}
\toprule
\textbf{Rubro} & \textbf{Plan original (S/)} & \textbf{Valor real (S/)} & \textbf{Veredicto} \\
\midrule
Dispositivo de prueba (laptop/tablet 16 GB) & 2\,500,00 & 2\,299–2\,599 & \textbf{Confirmado} \\
Almacenamiento externo 1 TB & 250,00 & 480–570 (SSD) / 280 (HDD) & \textbf{Subestimado} \\
Cómputo en la nube & 900,00 & 0–200 (niveles gratuitos) & \textbf{Sobreestimado} \\
Bibliografía y bases de datos & 300,00 & — & Razonable \\
Impresiones y empastados & 200,00 & — & Razonable \\
Transporte a zonas rurales (5 visitas) & 750,00 & — & Razonable \\
Viáticos para trabajo de campo (5 días) & 500,00 & — & Razonable \\
Material para talleres EIB & 200,00 & — & Razonable \\
Consulta con especialistas Quechua Collao & 450,00 & — & Razonable \\
Validación lingüística por hablantes nativos & 500,00 & — & Razonable \\
\midrule
\textbf{Subtotal} & 6\,550,00 & \textbf{6\,179,00} & \\
\textbf{Contingencia (10\%)} & 655,00 & \textbf{618,00} & \\
\midrule
\textbf{TOTAL} & \textbf{7\,205,00} & \textbf{$\approx$ 6\,800,00} & \\
\bottomrule
\end{tabular}
\end{table}

\textit{Nota:} los precios corresponden a cotizaciones en retailers peruanos (Hiraoka, Falabella, Promart) a junio de 2026, con tipo de cambio BCRP/SBS de $\approx$~S/ 3,40 por USD. Se recomienda \textbf{reasignar el excedente de la partida de nube} ($\approx$ S/ 700) a la corrección del almacenamiento y al refuerzo de las sesiones de validación lingüística.

\subsection{Factibilidad operativa y gestión del alcance}

Operativamente, el proyecto es más viable que una tesis dependiente de sistemas empresariales, porque el artefacto propuesto es un \textbf{prototipo autónomo} que no requiere acceso a código fuente, APIs internas o bases de datos corporativas. Sin embargo, mantiene dos dependencias operativas serias: la coordinación con docentes o especialistas EIB para la validación pedagógica, y el acceso a materiales curriculares en escenarios de prueba con pertinencia lingüística responsable.

\begin{table}[H]
\centering
\caption{Factores de riesgo operativo y estrategias de mitigación}
\label{tab:factibilidad_operativa}
\small
\begin{tabular}{@{}p{5.5cm}p{8.5cm}@{}}
\toprule
\textbf{Factor de riesgo} & \textbf{Estrategia de mitigación / delimitación} \\
\midrule
Acceso a validadores EIB (DRE/UGEL/AMLQ) & Iniciar gestiones institucionales desde el arranque del proyecto; tener plan B de validación remota con panel reducido de especialistas. \\
\midrule
«Agujero negro» del desarrollo UI/UX & Limitar la interfaz a un mínimo funcional (CLI o web local tipo Streamlit); concentrar el esfuerzo en el \textit{pipeline} SLM+RAG y su evaluación. El valor académico no está en la interfaz, sino en el \textit{pipeline} \citep{Hevia2025}. \\
\midrule
Indefinición del criterio de parada & Fijar umbrales de éxito \textit{ex ante}: (i) técnicos — latencia, RAM $\leq$ capacidad del equipo objetivo, tok/s, estabilidad \textit{offline}; (ii) lingüísticos — chrF++ sobre corpus de prueba Collao–español; (iii) pedagógicos — rúbricas de docentes EIB. \\
\midrule
Sobreextensión del corpus y alcance & Acotar el dominio curricular a áreas de Comunicación y Ciencia y Tecnología para el primer ciclo de validación; no intentar cubrir toda la EIB en una sola iteración. \\
\bottomrule
\end{tabular}
\end{table}

\subsection{Semáforo de riesgo tecnológico}

\begin{table}[H]
\centering
\caption{Semáforo de riesgo por dimensión}
\label{tab:semaforo}
\small
\begin{tabular}{@{}p{5cm}cp{8cm}@{}}
\toprule
\textbf{Dimensión} & \textbf{Nivel} & \textbf{Justificación} \\
\midrule
Técnica (hardware / inferencia / cuantización) & \textbf{VERDE} & Modelos y herramientas probados; inferencia en CPU validada; cuantización en nivel gratuito. \\
\midrule
Stack / tooling & \textbf{VERDE} & Maduro, abierto, documentado; alerta menor en tokenización del quechua. \\
\midrule
Económica / nube & \textbf{VERDE} & Costos cubribles; nube gratuita (Kaggle/Colab) cubre tareas \textit{one-off}; presupuesto corregido holgado. \\
\midrule
Datos (corpus Collao) & \textbf{AMARILLO} & Corpus base disponible \citep{Orcotoma2025} pero pequeño, sin limpieza \textit{post-OCR}; variantes dialectales (Chanka vs.\ Collao) exigen validación. \\
\midrule
Operativo (acceso EIB / alcance) & \textbf{AMARILLO} & Depende de trámites con DRE/UGEL/AMLQ y de disciplina de alcance para evitar \textit{scope creep}. \\
\midrule
\textbf{GLOBAL} & \textbf{AMARILLO} & \textbf{Factible con gestión activa de los riesgos de datos y operativos.} \\
\bottomrule
\end{tabular}
\end{table}

El proyecto pasa a \textbf{VERDE} si se consolidan estas condiciones: usar un modelo entre 3B y 7B cuantizado; mantener el sistema estrictamente textual y \textit{offline}; limitar la base de conocimiento a áreas curriculares concretas; fijar desde el inicio un corpus validado y un comité de revisión lingüística; y usar la nube solo como apoyo puntual, no como dependencia estructural \citep{Walusimbi2026, Xu2024}.

El proyecto pasa a \textbf{ROJO} si deriva hacia: entrenar o afinar intensivamente el modelo base como parte central de la tesis; incorporar voz, app móvil completa y despliegue institucional simultáneamente; o depender de una API externa para la generación principal.

\subsection{Conclusiones y recomendaciones}

La factibilidad es un espectro, no una decisión binaria, y este proyecto se ubica en el \textbf{lado factible}. La decisión de no entrenar desde cero, sino cuantizar y orquestar un SLM preexistente con RAG local, reduce drásticamente las exigencias de cómputo y costo, haciéndolo viable bajo cualquier escenario de hardware disponible para el equipo de tesis. Las recomendaciones concretas son las siguientes:

\begin{enumerate}
    \item \textbf{Modelo base:} adoptar \textbf{Qwen2.5-3B} como configuración por defecto (equipos de 8 GB) y promover el 7B si el equipo objetivo sostiene $\geq$ 8 tok/s con $\leq$ 7 GB de RAM durante pruebas reales.
    \item \textbf{Presupuesto:} subir el almacenamiento a S/ 500 (SSD) o S/ 300 (HDD) y reducir la nube a S/ 100–200, reasignando el excedente a validación lingüística.
    \item \textbf{Corpus:} acotar a un dominio curricular y grado específicos para el primer ciclo de validación, priorizando curación manual sobre cantidad.
    \item \textbf{Criterio de parada:} fijar antes del inicio del desarrollo los umbrales de latencia, RAM, chrF++ y rúbrica pedagógica como condiciones de suficiencia del prototipo.
    \item \textbf{Gestión institucional:} iniciar de inmediato las gestiones con la DRE/UGEL de Cusco (y eventualmente Puno) y la AMLQ para asegurar acceso a validadores.
    \item \textbf{Interfaz:} limitar el desarrollo a lo mínimo funcional; el valor académico reside en el \textit{pipeline} y su evaluación, no en la experiencia de usuario.
\end{enumerate}
